% ═══════════════════════════════════════════════════════════════════
% ПРИЛОЖЕНИЯ A–D
% ═══════════════════════════════════════════════════════════════════
\appendix

\part*{Приложения}
\addcontentsline{toc}{part}{Приложения}

\section{Доказательства $\Ga$-тождеств}\label{app:gamma}

\subsection{Формула отражения}

\begin{proposition}[интеграл Эйлера]\label{prop:euler}
Для $0<\re z<1$
\[
  \int_0^\infty \frac{u^{z-1}}{1+u}\,du \;=\; \frac{\pi}{\sin \pi z}.
\]
\end{proposition}

\begin{proof}
Рассмотрим «ключевой» контур вокруг разреза $[0,+\infty)$: большая
окружность радиуса $R$, две стороны разреза и малая окружность радиуса
$\rho$ вокруг нуля. Функция $f(u)=\frac{(-u)^{z-1}}{1+u}$ голоморфна
на $\CC\setminus[0,+\infty)$ при выборе ветви $(-u)^{z-1}$ с
аргументом в $(-\pi,\pi)$, и её единственный полюс $u=-1$ лежит внутри
контура с вычетом $\operatorname{Res}_{u=-1}f=(1)^{z-1}=1$. Интеграл
по внешней и внутренней окружностям стремится к нулю при $R\to\infty$,
$\rho\to0$ (поведение $u^{z-2}$ и $u^{z-1}$). Верхняя сторона разреза
($\arg(-u)\to-\pi$) проходится от $0$ к $\infty$ и даёт
$e^{-i\pi(z-1)}=e^{i\pi(1-z)}$ на интеграл; нижняя
($\arg(-u)\to+\pi$, обход от $\infty$ к $0$) даёт
$-e^{i\pi(z-1)}=-e^{-i\pi(1-z)}$. Итого по основной теореме о вычетах
\[
  \bigl(e^{i\pi(1-z)}-e^{-i\pi(1-z)}\bigr)
  \int_0^\infty\frac{x^{z-1}}{1+x}\,dx \;=\;2\pi i,
\]
то есть $\bigl(2i\sin\pi(1-z)\bigr)
\int_0^\infty\frac{x^{z-1}}{1+x}dx=2\pi i$; поскольку
$\sin\pi(1-z)=\sin\pi z$, получаем утверждение. Интеграл сходится при
$0<\re z<1$ (поведение $x^{z-1}$ при $x\to0$ и $x^{z-2}$ при
$x\to\infty$).
\end{proof}

\begin{proposition}[формула отражения]\label{prop:reflection}
Для $0<z<1$ имеем $\Ga(z)\Ga(1-z)=\frac{\pi}{\sin\pi z}$;
аналитическим продолжением тождество распространяется на все
$z\in\CC\setminus\ZZ$.
\end{proposition}

\begin{proof}
Замена $t=\frac{u}{1+u}$ (обратная $u=\frac{t}{1-t}$,
$du=\frac{dt}{(1-t)^2}$) переводит интеграл
предложения~\ref{prop:euler} в бета-интеграл:
\[
  \int_0^\infty\frac{u^{z-1}}{1+u}\,du
  \;=\;\int_0^1 t^{z-1}(1-t)^{-z}\,dt
  \;=\;\Bfun(z,1-z)
  \;=\;\frac{\Ga(z)\Ga(1-z)}{\Ga(1)}
  \;=\;\Ga(z)\Ga(1-z),
\]
где использовано представление Эйлера
$\Bfun(p,q)=\int_0^1t^{p-1}(1-t)^{q-1}dt=\Ga(p)\Ga(q)/\Ga(p+q)$ при
$\re p,\re q>0$ и $\Ga(1)=1$. Сопоставление с
предложением~\ref{prop:euler} даёт формулу при $0<z<1$; обе части ---
мероморфные функции, совпадающие на полосе, значит, совпадают всюду
вне полюсов.
\end{proof}

\subsection{Формула умножения Гаусса}

\begin{proposition}[формула умножения]\label{prop:mult}
Для целых $m\ge2$ и $z$ с $mz\notin\ZZ_{\le0}$
\[
  \prod_{k=0}^{m-1}\Ga\Bigl(z+\frac{k}{m}\Bigr)
  \;=\;(2\pi)^{\frac{m-1}{2}}\;m^{\frac12-mz}\;\Ga(mz).
\]
\end{proposition}

\begin{proof}
Обозначим левую часть $G_m(z)$, правую ---
$H_m(z)=(2\pi)^{\frac{m-1}{2}}\,m^{\frac12-mz}\,\Ga(mz)$.
Рекуррентность $\Ga(w+1)=w\,\Ga(w)$ даёт
\[
  \frac{G_m(z+1)}{G_m(z)}
  \;=\;\prod_{k=0}^{m-1}\Bigl(z+\frac{k}{m}\Bigr)
  \;=\;m^{-m}\,\frac{\Ga(mz+m)}{\Ga(mz)},
\]
ибо $\prod_{k=0}^{m-1}(mz+k)=\Ga(mz+m)/\Ga(mz)$. Та же рекуррентность
даёт
\[
  \frac{H_m(z+1)}{H_m(z)}
  \;=\;m^{-m}\,\frac{\Ga(mz+m)}{\Ga(mz)}
\]
(множитель $m^{\frac12-mz}$ при сдвиге умножается на $m^{-m}$).
Следовательно, отношение $Q_m=G_m/H_m$ есть $1$-периодическая
голоморфная функция на $\CC\setminus(-\infty,0]$. Формула Стирлинга
даёт при $\re z\to+\infty$ в вертикальной полосе конечной ширины
$Q_m(z)\to1$: асимптотика Стирлинга по каждому сомножителю даёт
\[
  \prod_{k=0}^{m-1}\Ga\Bigl(z+\frac{k}{m}\Bigr)
  \;\sim\;(2\pi)^{\frac{m-1}{2}}\,m^{\frac12-mz}\,\Ga(mz)
  \qquad(\re z\to+\infty),
\]
и числитель со знаменателем сокращаются. Периодичность переносит
предел на любую вертикаль полосы, откуда $Q_m\equiv1$ (ограниченная
периодическая функция с нулевым множеством предельных значений
постоянна). Приведённый аргумент самодостаточен для всех $m$ сразу.
\end{proof}

\begin{remark}
Для нужд стенда достаточно следствий: рекуррентности
$\Ga(z+1)=z\Ga(z)$ (использована в лемме~\ref{lem:relations},
пункт~2), отражения и умножения. Все три выведены выше из
интегральных представлений и основной теоремы о вычетах --- внешние
результаты не привлекаются.
\end{remark}

\section{Перепись характеров по проводникам}\label{app:census}

\subsection{Формула подсчёта}

Пусть $d\mid N$ и $g_0=N/d$. По определению~\ref{def:conductor}
$h_d=\#\{(a,b)\in T_N:\ \gcd(N,a,b)=N/d\}$. Условие
$\gcd(N,a,b)=g_0$ равносильно тому, что $g_0\mid a$, $g_0\mid b$ и
$\gcd\bigl(\frac{a}{g_0},\frac{b}{g_0},d\bigr)=1$. Полагая
$\alpha=a/g_0$, $\beta=b/g_0$ и используя
$\alpha+\beta\le\frac{N-1}{g_0}=d-\frac1{g_0}$ (что для целых
$\alpha+\beta$ означает $\alpha+\beta\le d-1$), получаем
\[
  h_d\;=\;\#\Bigl\{(\alpha,\beta):\ \alpha,\beta\ge1,\
  \alpha+\beta\le d-1,\ \gcd(\alpha,\beta,d)=1\Bigr\}.
\]
Через функцию Мёбиуса: обозначим
$c(k)=\#\{(\alpha,\beta):\ \alpha,\beta\ge1,\ \alpha+\beta\le k-1\}
=\frac{(k-1)(k-2)}{2}$; тогда
\begin{equation}\label{eq:mobius}
  h_d\;=\;\sum_{m\mid d}\mu(m)\;c\!\Bigl(\frac{d}{m}\Bigr)
  \;=\;\sum_{m\mid d}\mu(m)\,
  \frac{\bigl(\frac{d}{m}-1\bigr)\bigl(\frac{d}{m}-2\bigr)}{2}.
\end{equation}

\subsection{Таблицы для $N=15$ и $N=30$}

Для $N=15$ (проводники $3,5,15$):
\begin{center}
\begin{tabular}{@{}lccc@{}}
\toprule
$d$ & Вычисление по \eqref{eq:mobius} & $h_d$ & Контроль \\
\midrule
$3$  & $c(3)-c(1)=1-0$ & $1$  & \\
$5$  & $c(5)-c(1)=6-0$ & $6$  & \\
$15$ & $c(15)-c(5)-c(3)+c(1)=91-6-1+0$ & $84$ & \\
\midrule
\multicolumn{3}{@{}l}{Итого: $1+6+84$} & $=\;91=g(15)$ \\
\bottomrule
\end{tabular}
\end{center}

Для $N=30$ (проводники $3,5,6,10,15,30$):
\begin{center}
\small
\begin{tabular}{@{}l p{6.6cm} c c@{}}
\toprule
$d$ & Вычисление по \eqref{eq:mobius} & $h_d$ & Контроль \\
\midrule
$3$  & $c(3)-c(1)$ & $1$ & \\
$5$  & $c(5)-c(1)$ & $6$ & \\
$6$  & $c(6)-c(3)-c(2)+c(1)=10-1-0+0$ & $9$ & \\
$10$ & $c(10)-c(5)-c(2)+c(1)=36-6-0+0$ & $30$ & \\
$15$ & $c(15)-c(5)-c(3)+c(1)=91-6-1+0$ & $84$ & \\
$30$ & $c(30)-c(15)-c(10)-c(6)+c(5)+c(3)+c(2)-c(1)$ & $276$ & \\
     & $=\,406-91-36-10+6+1+0-0$ & & \\
\midrule
\multicolumn{3}{@{}l}{Итого: $1+6+9+30+84+276$} & $=\;406=g(30)$ \\
\bottomrule
\end{tabular}
\end{center}

Проверка $h_{30}$ по формуле включений-исключений: множители
$30=2\cdot3\cdot5$, знаки $\mu(m)$: $+1,-1,-1,-1,+1,+1,+1,-1$ для
$m=1,2,3,5,6,10,15,30$. В скрипте-протоколе обе таблицы получены
прямым перебором пар $(a,b)$ и сверены с формулой~\eqref{eq:mobius}
точной целочисленной арифметикой (V1).

\section{Таблицы данных сертификатов}\label{app:data}

\subsection{Семейство фаз и торможений тора}

\begin{center}
\small
\begin{tabular}{@{}c c c c c@{}}
\toprule
$N$ & $\delta=\pi/N$ & $\gamma=\delta^4$ ($k=1$) &
$\delta_{\mathrm{eff}}$ ($k=1$) & $\gamma$ ($k=22$) \\
\midrule
2 & 1.5707963 & 6.0880682 & 9.5631151 & 0.2767304 \\
3 & 1.0471976 & 1.2025814 & 1.2593403 & 0.0546628 \\
4 & 0.7853982 & 0.3805043 & 0.2988473 & 0.0172956 \\
5 & 0.6283185 & 0.1558545 & 0.0979263 & 0.0070843 \\
6 & 0.5235988 & 0.0751642 & 0.0393464 & 0.0034166 \\
7 & 0.4487990 & 0.0404989 & 0.0181767 & 0.0018409 \\
8 & 0.3926991 & 0.0237754 & 0.0093371 & 0.0010807 \\
\bottomrule
\end{tabular}
\end{center}

\subsection{Ключевые константы стенда}

\begin{center}
\small
\begin{tabular}{@{}l l l@{}}
\toprule
Величина & Значение & Источник \\
\midrule
$b_{Ch}(7)$ & $\approx 0.7526$ & $1-\cos(2\pi/7)$ \\
$b_{Ch}(9)$ & $\approx 0.1172$ & $1-\cos(2\pi/9)$ \\
$b_{Ch}(11)$ & $\approx 0.0823$ & $1-\cos(2\pi/11)$ \\
$b_{Ch}(15)$ & $\approx 0.086454$ & радикал~\eqref{eq:b15} \\
$b_{Ch}(30)$ & $\approx 0.021854$ & радикал~\eqref{eq:b30} \\
$\disc$ стенда & $1265625=1125^2$ & $3^4\cdot5^6$ \\
SNF-тип & $(1,1,5,5,15,15,15,15)$ & теорема~\ref{th:H1} \\
$Q_{\text{стенд}}$ & $480$ & теорема~\ref{th:H56} \\
$j$ Клейна & $-3375=-15^3$ & $\Delta=-7^3$ \\
$\tau$ Клейна & $(1+\sqrt{-7})/2$ & $\im\tau=\sqrt7/2$ \\
\bottomrule
\end{tabular}
\end{center}

\section{Воспроизводимость и верификация}\label{app:repro}

\subsection{Лаборатория \texttt{laboratory.py}}

Единый файл \texttt{laboratory.py} в корне репозитория реализует
полную вычислительную лабораторию программы:
\begin{itemize}[leftmargin=1.6em, itemsep=0.1em]
\item меню запуска тестов (полный протокол V1--V9, стенды, серт.);
\item переключение языков интерфейса (русский/английский);
\item интерактивные параметры расчётов ($dps$, выборки, уровни $N$);
\item конструктор собственных экспериментов (новые лабораторные
  опыты с произвольными параметрами);
\item генерация отчётов в отдельную папку \texttt{reports/} (JSON +
  текст);
\item плиточные диаграммы в высоком разрешении $600$~dpi;
\item мультиязычная верификация: Python, Julia, Fortran, C, Rust;
\item интеграция Lean: точные целочисленные тождества проверяются
  ядром Lean~4.
\end{itemize}

\subsection{Карта интеграции литературы}

Сертификат H зафиксировал карту интеграции программы с
классическими источниками: Коула--Селберг (PNAS 35, 1949) ---
$\Ga$-слой периодов; Гросс (LNM 588, 1980) --- нормировка периодов
CM; Гросс---Рорлих (Invent.\ 44, 1978) --- якобианы кривых Ферма;
Вейль (1949) --- характерное разложение якобиана Ферма;
PSLQ Фергюсона---Бейли (1994) --- независимое восстановление
экспонент в служебных проверках; LLL (1982) --- фазы в поле
классов; Ленстра (2016) --- интерфейс масштабирования на
произвольные многообразия. Каждый элемент карты либо используется в
служебной роли с независимой сверкой, либо подтверждён замкнутой
формой основного текста.
